% Encodage des caractères et langue du document
\usepackage{setspace}
% \usepackage[utf8]{inputenc}

% ---- fonts
%\usepackage{lmodern} % latin modern font
\usepackage{newtxtext,lmodern} % Times font for text and lmodern for maths
%\usepackage{fourier} % Utopia font
\usepackage[english]{babel}
\usepackage[T1]{fontenc} % mandatory for hyphenation
\usepackage{textgreek}
\usepackage{bbm}
\usepackage{ifthen}
\usepackage[strict]{changepage}
\usepackage{titletoc}
\usepackage{supertabular}
\usepackage[para]{threeparttable}
\usepackage{lscape}
% \usepackage[latin1]{inputenc}
\usepackage[cyr]{aeguill}
\usepackage{xspace}
\usepackage[version=4]{mhchem}  % for chemical species
\usepackage[most]{tcolorbox}    % better boxes

\setcounter{secnumdepth}{3} % Pour que les subsubsections ne soient pas numérotées
\usepackage{fixltx2e}
\usepackage{hyperref}
\usepackage{eso-pic} % for footer image

\usepackage[final]{pdfpages}

%\includeonly{chap3/chap3_main}  % in order to compile only specific chapters

%%%%%%%%%% Gestions des marges %%%%%%%%%% 
\usepackage[twoside]{geometry} % Si on a besoin d'une configuration plus précise des marges
\geometry{a4paper,                % format de papier
% Définition des marges :
  left= 2.8cm,right = 2cm,  % marge intérieure extérieure à la page
  top = 3cm,bottom = 3cm,
% En-tête et pied de page :
  headheight=6mm,         % espace réservé à l'en-tête dans la marge top
  %headsep=3mm,            % espace entre le corps et l'en-tête
  %footskip=9mm            % espace entre le corps et le pied de page
  marginparwidth = 16mm,
}

\raggedbottom
\reversemarginpar
%\usepackage{showframe} % pour afficher les traits des marges

%%%%%%%%%%%%%%%%%%%%%%%%%%%%%%%%%%%%%%%%%%%%%%%%%%%%%%%%%%%%%%%%%%%%
%%%%%%%%%%% Gestion maths %%%%%%%%%% 
\usepackage{amsmath,amssymb,amsfonts,amsthm}
\usepackage{mathtools} % version modifiée de amsmath, ajoute des symboles, etc.
\usepackage{mathrsfs}% pour rajouter un format de lettres façon calligraphie en math mode.
\DeclareMathOperator{\sinc}{sinc}
\DeclareMathOperator{\e}{e}

\usepackage[locale = UK]{siunitx} % Pour gérer les unités
\sisetup{
  % Numbers
  group-separator          = {\,},     % thin space for thousands
  group-minimum-digits     = 4,        % start grouping from 10 000
  output-decimal-marker    = {.},      % decimal point
  exponent-product         = \times,  % scientific notation with ×
  output-exponent-marker   = \mathrm{e}, % optional: 1e5 instead of ×10^5

  % Units
  per-mode                 = power,   % use · instead of  
  inter-unit-product       = \ensuremath{{}\cdot{}}, % keep ⋅ 
  detect-weight            = true,     % bold in math bold
  detect-family            = true,     % italics in section headings

  % Ranges and lists
  range-phrase             = {\--}, % en dash in ranges
  range-units              = single,     % do not repeat units
  list-separator           = {, },       % 10 m, 20 m, 30 m
  list-final-separator     = {, and }    % 10 m, 20 m, and 30 m
}
%\sisetup{separate-uncertainty=true,multi-part-units=single} % pour faire des incertitudes en écrivant \SI{valeur(incertitude)}{unité}
\DeclareSIUnit\vitesse{\meter\per\second}
\usepackage{eurosym}
\DeclareSIUnit{\octet}{o}

\def\Oz{O\textsubscript{3}}
\def\Ammo{NH\textsubscript{3}}

%%%%%%%%%%%%%%%%%%%%%%%%%%%%%%%%%%%%%%%%%%%%%%%%%%%%%%%%%%%%%%%%%%%%
%%%%%%%%%%%  Graphics / Table / List %%%%%%%%%%%
\usepackage{graphicx,array,tikz,multirow,adjustbox}
\usepackage{caption,subcaption} % permet de faire des subfigures (remplace le package subfig)
\usepackage{svg,float}
\usepackage{booktabs,paralist}
\newcolumntype{x}[1]{>{\centering\arraybackslash\hspace{0pt}}p{#1}}
\usepackage[section]{placeins}
\usepackage{hanging}
\usepackage{sidecap}
\sidecaptionvpos{figure}{t}

\captionsetup{font=small, width=\textwidth}
\usepackage{supertabular}



%%%%%%%%%%%%%%%%%%%%%%%%%%%%%%%%%%%%%%%%%%%%%%%%%%%%%%%%%%%%%%%%%%%%
%%%%%%%%%%% Header / Foot %%%%%%%%%%% 
\usepackage{fancyhdr,emptypage} % garantit que les pages blanches avant les débuts de chapitres soient vraiment blanches (pas d'en-tête ni de pied de page)
\let\cleardoublepage\clearpage
 
\fancypagestyle{plain}{ %% Page chapitre, toc ...
    %\fancyhead{}\fancyfoot[C]{\thepage}
    \renewcommand{\headrulewidth}{0pt}
    \renewcommand{\footrulewidth}{0pt}
}

%%%%%%%%%%% Page normale
\pagestyle{fancy}
    %\renewcommand{\chaptermark}[1]{\markboth{\color{chapcolor}\thechapter.\ #1}{}}
    %\renewcommand{\chaptermark}[1]{\markboth{\chaptername \ \thechapter.\ \textsc{#1}}{}} % sert à personnaliser l'affichage de \leftmark (ici : le mot "Chapitre", le numéro, un point, et le titre du chapitre, sans écrire en majuscules)
    % \renewcommand{\sectionmark}[1]{\markright{\thesection.\ #1}} % sert à personnaliser l'affichage de \rightmark (ici : le numéro et le titre de la section en cours, sans écrire en majuscules)
    % \renewcommand{\sectionmark}[1]{\markright{\uppercase\expandafter{\romannumeral\thecounterchapter{}\relax}.\arabic{section}.\ #1}} % JULIE : pour mettre en haut à droite le numéro du chapitre puis la section
    \fancyhf{} % assure que les entête et pieds de page sont vides au départ

    % ANDERSON : header avec chapitre page de gauche et section page de droite
    %% --------- Définition chaptermark et sectionmark pour les entêtes ---------    
    \renewcommand{\chaptermark}[1]{\markboth{\chaptername \ \thechapter.\ \textsc{#1}}{}}
    \renewcommand{\sectionmark}[1]{\markright{\thesection.\ #1}}
    %\renewcommand{\sectionmark}[1]{\markright{\textit{#1}}}
    %%(ATTENTION, commandes délicate)
    %\renewcommand{\thechapter}{\Roman{chapter}}
    %\fancyhead[LO]{\thesection}
    %\fancyhead[RO]{\rightmark}
    %\fancyhead[LE]{\textsc{Chapter}\ \thechapter}
    %\fancyhead[RE]{\leftmark}

    %\fancyhead[LE]{\textsc{Chapter}\ \thechapter}
    %\fancyhead[RE]{\selectfont\nouppercase{\leftmark}}
    \fancyhead[LE]{\selectfont\leftmark}
    \fancyhead[RO]{\selectfont\rightmark} 
    \fancyfoot[LE,RO]{\thepage} % pied de page avec le numéro à droite et à gauche

% Explications :
% L = left, R = right, C = center, E = even pages, O = odd pages
%\leftmark : adds name and number of the current top-level structure (for example, Chapter for reports and books classes; Section for articles ) in uppercase letters.
%\rightmark : adds name and number of the current next to top-level structure (Section for reports and books; Subsection for articles) in uppercase letters.

%%%%%%%%%%% Personnaliser les premières pages des chapitres

\usepackage[Lenny]{fncychap}
\ChNameVar{\fontsize{25}{25}\selectfont\textcolor{chapcolor}}
\ChRuleWidth{0pt}
\ChNumVar{\fontsize{60}{62}\selectfont\textcolor{chapcolor}}

\makeatletter
\ChTitleVar{\Huge\rm}
\renewcommand{\DOCH}{%
\setlength{\fboxrule}{\RW} % Let fbox lines be controlled by
\fbox{\CNV\FmN{Chapter}\space \CNoV\thechapter}\par\nobreak
\vskip 20\p@}
\renewcommand{\DOTIS}[1]{%
\CTV\bfseries\FmTi{#1}\par\nobreak
\vskip 20\p@}
\makeatother

\usepackage[explicit]{titlesec}
\usepackage{lmodern}
\usepackage{lipsum}
\usepackage{color}

\newlength\chapnumb
\setlength\chapnumb{3cm}
 
\titleformat{\chapter}[block]{
  \normalfont}{}{0pt} {
    \parbox[b]{\chapnumb}{
      \fontsize{100}{100}\selectfont\textcolor{chapcolor}\thechapter}
      \parbox[b]{\dimexpr\textwidth-\chapnumb\relax}{
        \raggedleft
        \hfill{\color{chapcolor} \Huge \textbf{#1}}\\
        \textcolor{chapcolor}{\rule{\dimexpr\textwidth-\chapnumb\relax}{0.4pt}}
  }
}
 
\titleformat{name=\chapter,numberless}[block]
{\normalfont}{}{0pt}
{\parbox[b]{\chapnumb}{%
   \mbox{}}%
  \parbox[b]{\dimexpr\textwidth-\chapnumb\relax}{%
    \raggedleft%
    \hfill{\LARGE#1}\\
    \rule{\dimexpr\textwidth-\chapnumb\relax}{0.4pt}}}

\renewcommand{\thechapter}{\Roman{chapter}}

\renewcommand{\thesection}{\arabic{section}}

% Pour la table des matières
% \usepackage{minitoc}	
% \usepackage[french,nohints,tight]{minitoc}		% Mini table des matières, en français
%\setcounter{minitocdepth}{6} % Mini-toc détaillées (sections/sous-sections)
\setcounter{tocdepth}{2} % 
%\setlength{\mtcindent}{-1em} % décalage des minitoc à gauche
%\dominitoc

\usepackage[nottoc]{tocbibind} % pour que la bibliographie apparaisse dans la table des matières (avec l'option pour que la table des matières elle-même n'apparaisse pas dans la table des matières).
% \usepackage{tocloft}% pour pouvoir modifier les tailles d'espacement dans la table des matières
\usepackage[titles]{tocloft}

%%%%%%%%%%%%%%%%%%%%%%%%%%%%%%%%%%%%%%%%%%%%%%%%%%%%%%%%%%%%%%%%%%%%
%%%%%%%%%%% Divers %%%%%%%%%%% 
\usepackage{textcomp} % rajoute des symboles 
\usepackage{xcolor} % pour ajouter de la couleur (si besoin)
\usepackage{epigraph} % pour rajouter des citations en début de chapitre  \epigraph{Citation}}{Auteur}
\usepackage{titling}
\usepackage{lipsum} 
\usepackage{csquotes} % added 07/09/21 : \begin{displayquote}

\usepackage{xspace}
\usepackage{afterpage}
\renewcommand{\baselinestretch}{1.2} % interligne

\usepackage[textsize=footnotesize]{todonotes}

%%%%%%%%%%%%%%%%%%%%%%%%%%%%%%%%%%%%%%%%%%%%%%%%%%%%%%%%%%%%%%%%%%%%
%%%%%%%%%%% Links ref  %%%%%%%%%%% 
\usepackage{bookmark}
\usepackage{acronym}
\usepackage[nameinlink,english]{cleveref} % noabbrev
\Crefname{figure}{Fig.}{Figs.} % traduction des références aux figures/tables/équations
    \crefname{figure}{Fig.}{Figs.} % {ref type}{abbrev for one}{abbrev for multiple}
\Crefname{equation}{Eq.}{Eqs.}
\crefname{equation}{eq.}{eqs.}
\Crefname{table}{Table}{Tables}
\crefname{table}{table}{tables}
\crefname{section}{Sect.}{Sects.}
\crefname{chapter}{Chap.}{Chaps.}

% Configuration de hyperref
\definecolor{color_ref}{rgb}{0.18, 0.31, 0.31} % couleur cite
%\definecolor{color_link}{RGB}{36, 56, 141}
\definecolor{color_link}{RGB}{18, 69, 77}
\definecolor{curcolor}{RGB}{113,127,184} 
\definecolor{chapcolor}{RGB}{10, 77, 41} % couleur des titres de chapitres
\definecolor{boxcolor}{HTML}{FFF5D1} % couleur des box
\definecolor{OQboxcolor}{HTML}{FFF5D1} % couleur OQ box
\definecolor{KRboxcolor}{HTML}{E3F1EC} % couleur KR box
\definecolor{persoblue}{HTML}{CBFCFF}
\definecolor{persolavender}{HTML}{967BB6}
\definecolor{persogreen}{HTML}{6B8E23}
\definecolor{persored}{HTML}{CD5C5C}
\definecolor{persoorange}{HTML}{CA6F1E}
\definecolor{persogrey}{HTML}{7EA9A5}
\definecolor{persopink}{HTML}{B26A98}
\definecolor{persobluedark}{HTML}{1B2892}


\hypersetup{
	colorlinks=true, % colore les liens au lieu de les encadrer
	pdfstartview=FitV, % ouvre le PDF de façon à ce qu'il prenne la taille verticale de l'écran
	urlcolor=color_link, % choix de la couleur des liens URL
	linkcolor= color_link, % choix de la couleur des liens internes (table des matières, etc.)
	citecolor=color_ref % choix de la couleur des liens de citations
}

%%%%%%%%%%%%%%%%%%%%%%%%%%%%%%%%%%%%%%%%%%%%%%%%%%%%%%%%%%%%%%%%%%%%
%%%%%%%%%%% Bibliography ref  %%%%%%%%%%%

\usepackage[hyperref=true,natbib=true,
			backref=true,date=year,
			backend=biber,
			url=false,doi=true,isbn=false,%
			minbibnames=1, % nb min authors in biblio
			maxbibnames=6, % nb max authors in biblio
			maxcitenames=1, % nb min authors as textual
      mincitenames=1, % nb min authors as textual
			% maxalphanames=1, %nb author ref
			style=authoryear,%authoryear, numeric ,  alphabetic
			sorting=nyt,uniquename=false,uniquelist=false]{biblatex}
\renewcommand*{\bibfont}{\footnotesize}
\setlength\bibitemsep{\itemsep}
\renewbibmacro{in:}{} % remove In 

% \DeclareUnicodeCharacter{2212}{-}

%\renewcommand*{\labelalphaothers}{}
\DeclareLabelalphaTemplate{
  \labelelement{
    \field[final]{shorthand}
    \field{labelname}
    \field{label}
  }
  \labelelement{\literal{,\addhighpenspace}}
  \labelelement{\field{year}}
}

\AtEveryBibitem{%
    \clearfield{note} % Remove note
    \clearlist{language} % Remove doi
}


\usepackage{pdfpages}

\usepackage{colortbl}


